%% copyright --- strona po stronie tytułowej.
%%
%% \bookedition i \bookyear są przypięte w preamble.tex, a \pinnedon jest
%% napisem w lang/<lang>.tex. Żadnej z tych trzech wartości nie wpisuje się
%% tu ręcznie.

\thispagestyle{empty}
\vspace*{\fill}

{\small\color{mgrey}\raggedright\setlength{\parindent}{0pt}%
\emph{Python dla inżynierów .NET}. Wydanie \bookedition.

\medskip
Copyright \textcopyright{} \bookyear{} Konrad Cinkusz. Wszelkie prawa
zastrzeżone poza zakresem udzielonych tu licencji.

\medskip
Tekst jest na licencji CC BY-NC-SA 4.0. Kod \dash{} każdy plik w katalogu
\texttt{code/} i każdy listing, który książka z niego drukuje \dash{} jest na
licencji MIT, więc można go wkleić do pracy komercyjnej bez namysłu.

\medskip
Wydane własnym nakładem autora. Nie ma wydawcy, dystrybutora ani numeru
ISBN, więc tę książkę można czytać i kopiować na powyższych licencjach, ale
nie może jeszcze trafić do księgarni ani do katalogu biblioteki.

\medskip
Wydana w dwóch językach, angielskim i polskim, z jednego repozytorium.
Dostępna także po angielsku jako \emph{Python for .NET Engineers}.

\medskip
Wersje przypięto: \pinnedon. Każdy listing w tej książce uruchomiono na
wersjach ze strony tytułowej, a build repozytorium uruchamia je ponownie przy
każdej zmianie. To inna data niż rok w nocie copyright powyżej i znaczy co
innego.

\medskip
Python jest znakiem towarowym Python Software Foundation. .NET i C\# są
znakami towarowymi Microsoft Corporation. Ta książka nie jest z nimi
powiązana ani przez nie autoryzowana.
\par}

\clearpage
